\cotefile{83}
\scanpage{17}
\mspage{7\handwritten{5} \deleted{\unsure{bis}}}
% archive sheet number: 16

J'ai repris \inserted{\handwritten{tantôt}} l'expression ``psychodrame'', terme assez à la mode je \deleted{\unsure{croi}}
crois\handwritten{,} et qui ne manque pas de force\handwritten{.} \deleted{évocatrice.} Et il est vrai que scénario
et mise en scène de beaucoup de rêves\deleted{, \unsure{ceux surtout qui virent dans le cauchemar,}}
ne ménagent pas les sombres couleurs, aux tonalités \deleted{\unsure{dramati-}}
de \deleted{\unsure{dram}} ``drame'' antique \deleted{\unsure{dans la t}} repris dans la trame de notre vie à
nous, quand \handwritten{elles} ne virent à l'angoisse pure du cauchemar. Nous, acteurs
du drame, \inserted{y} ``marchons'' à fond, certes\handwritten{;} quitte\handwritten{,} au réveil\handwritten{,} \deleted{\unsure{de}} \inserted{à} se sentir tout
idiot et à\deleted{\unsure{x}} se dépêcher de penser à autre chose\handwritten{!} \deleted{\unsure{La}} \handwritten{C}e n'est pas nous, ça
c'est sûr, qui ``parlons'' cette langue étrange, la \deleted{``langue-parabole''}
``langue-images'' ou ``\uline{langue}-\uline{parabole}'', la ``\uline{langue}-\uline{vie}''. Il y\handwritten{\unsure{,}} a un Metteur
en Scène, un plus grand que nous, qui la manie comme \deleted{\unsure{sa}} \inserted{\handwritten{une}} langue maternelle\handwritten{,}
dont nous-mê\handwritten{\unsure{\^{}}}me et la substance même de notre vie \deleted{serions} \inserted{formerions} la chair des
mots. Il monte pêle-mêle drames, farces, idylles ou élégies -- mais même
là où le drame bat son plein, je vois \inserted{\handwritten{pourtant}} que \deleted{\unsure{Celui qui fait jouer fait le}} \inserted{le Narrateur invisible} \deleted{qui fait le récit, il}
garde un sourire en coin\handwritten{\unsure{.|!}} Il sait que la souffrance et la
mort s\handwritten{ont} des choses \inserted{\handwritten{toutes}} simples\handwritten{,} \deleted{et importantes, et} qui n'ont rien de ``dramatique''.
\deleted{\unsure{xx}} Le ``drame'', c'est les vagues que nous \inserted{\handwritten{aimons}} \deleted{\unsure{faisons}} fai\handwritten{re} autour de ces choses\handwritten{,}
histoire de \inserted{\handwritten{nous}} les rendre compliquées...

C'est de \deleted{me rappeller} \inserted{\handwritten{sentir}} ce ``sourire en coin'' du Narrateur, qui fait
que finalement, je \deleted{n'arrive pas \unsure{t}} n'arrive pas à \deleted{\unsure{être entièrement}} \inserted{me sentir} à l'aise
avec \deleted{\unsure{tout un} ``} \handwritten{c}e nom ``psychodrame''\handwritten{,} \deleted{qui ne manque pas} Il a beau avoir de
la gue\deleted{\unsure{l}}\unsure{u|l}le -- je préfère \deleted{\unsure{parler de}} \inserted{\handwritten{le nom serein de}} ``parabole''\handwritten{.} \deleted{aux connotations plus paisibles}
Une parabole \deleted{qui est} ``mise en scène'', c'est une chose entendue\handwritten{,} \deleted{Avec} \inserted{\handwritten{avec un ``scénario'' plus ou moins simple ou plus ou moins imbriqué.}}
\deleted{une} \handwritten{M}ise en scène souvent cocasse, toujours im\unsure{pré|pér}vue et toujours incisive
(sans nul souci, faut il croire, si elle sera comprise...). Elle n'est
limitée par aucune convention\handwritten{,} de style \handwritten{ou} de biènséance ($^{x}$), ni par \handwritten{aucune}
limitation\deleted{s} technique\deleted{s} pour les montages à grand spectacle -- des montages
à faire pâlir les réalisations les plus fantastiques du \handwritten{\unsure{magi}}cien-cinéaste
le plus inspiré\handwritten{!} Si tel ``magicien'' de l'image et \deleted{\illegible{1}} du son parfois nous enchante
ou nous touche, c'est pour avoir su écouter en lui-même, sûrement,
ce ma\handwritten{\unsure{î}}tre-Magicien aux moyens prodigieux, en oeuvre en chacun de nous, et
que si souvent nous dédaignons d'écouter. Il est vrai que ses spectacles
sont gratuits, et qu'il ne \handwritten{f}ait pas de publicité\handwritten{\unsure{.|!}}

\noindent\rule{.23\linewidth}{.3pt}
% slanted
\margin{\handwritten{\unsure{espacer|espace}}}

Ce qui \handwritten{s}ans doute déconcerte le plus dans la langue-paraboles, la
langue-images, c'est la \mbox{l\,i\,b\,e\,r\,t\,é}. On sent que c'est là son âme et
son essence\handwritten{,} \deleted{\unsure{qui}} \handwritten{Elle} a de quoi effrayer plus d'un : une liberté créatrice infinie.
Aucune règle à quoi se raccrocher, ni pour la ``parler'' (mais la
question ne se pose guère -- c'est \deleted{\unsure{quelqu}} un \deleted{\unsure{a}} \handwritten{A}utre en nous qui la parle\handwritten{..)\unsure{,}}
ni pour l'entendre et la saisir -- au vol ! Aucun pas ne semble prescrire
le suivant, ni être prescrit par celui qui l'a précédé -- \deleted{et pourtant,}
